% GENERATED FILE. Edit canonical lessons, then run npm run generate:books.
% canonical-source-hash: fnv1a64:67d2b56f776124b5
% canonical-lessons: LA-C482-amnis, LA-C482-torrens, LA-C482-confluens, LA-C482-fluo, LA-C482-transnato

\chapter{River, Torrent, Confluence, To flow, To swim across}
\label{ch:la-river-torrent-confluence-to-flow-to-swim-across}
\hlchaptermodality{482}

\begin{chapteropening}
\textbf{By the end of this chapter:} \emph{I can describe rivers: a broad river, a torrent, the place where rivers meet, water flowing, and swimming across.}

Complete the last lesson of chapter 482: I can describe rivers: a broad river, a torrent, the place where rivers meet, water flowing, and swimming across.
\end{chapteropening}

\section[amnis, amnis]{amnis, amnis — a river, a broad stream}

\label{lesson:LA-C482-amnis}



\begin{quote}
Before the new one: say the Latin for rough, then the Latin for to slide down.
\end{quote}

\subsection*{You'll want to know: amnis, amnis}
\textbf{amnis, amnis} — \textquotedblleft{}a river, a broad stream\textquotedblright{}.

\begin{quote}
Amnis magnus et altus per campōs it.
\end{quote}

\begin{quote}
\emph{A big, deep river runs through the plains.}
\end{quote}

Third declension, masculine: \emph{amnis, amnis m.} A big, deep river, a word more of poetry and history than \emph{flūmen}.

\begin{cousinweb}
\emph{amnis} is inherited from Old Latin \emph{abnis}, from PIE \emph{h\textsubscript{2}ep-} 'water' (EtymDB). Lewis \& Short also derive it from the Sanskrit word for water.
\end{cousinweb}

\begin{grammarlens}[title={What you've built}]
One more word for telling, explaining and describing at B1.
\end{grammarlens}

\subsection*{Your turn}
\begin{itemize}
\raggedright
  \item \emph{Say it:} \emph{amnis, amnis}
  \item \emph{Say it:} \emph{amnis, amnis}, once more
  \item \emph{Say it:} a sentence of your own with \emph{amnis, amnis}
  \item \emph{Recall:} say the Latin for rough, then the Latin for to slide down, then say \emph{amnis, amnis} again
\end{itemize}

\begin{tcolorbox}[breakable,colback=teal!4,colframe=teal!35!black,title={Before you move on}]
What does amnis, amnis mean? (\textquotedblleft{}A river, a broad stream\textquotedblright{}.) Say it once more.
\end{tcolorbox}
\section[torrēns, torrentis]{torrēns, torrentis — a torrent, a rushing stream}

\label{lesson:LA-C482-torrens}



\begin{quote}
Before the new one: say the Latin for to slide down, then the Latin for a river.
\end{quote}

\subsection*{You'll want to know: torrēns, torrentis}
\textbf{torrēns, torrentis} — \textquotedblleft{}a torrent, a rushing stream\textquotedblright{}.

\begin{quote}
Post pluviam torrēns ē monte currit.
\end{quote}

\begin{quote}
\emph{After the rain a torrent rushes down from the mountain.}
\end{quote}

Third declension, masculine: \emph{torrēns, torrentis m.} The present participle of \emph{torreō} 'to burn', used as a noun: a stream that is fierce. English \emph{torrent} comes from it.

\begin{cousinweb}
\emph{torrēns} is the present participle of \emph{torreō} 'to parch, to burn' (EtymDB; Lewis \& Short), used as a noun for a raging, boiling stream. English \emph{torrent} comes from it.
\end{cousinweb}

\begin{grammarlens}[title={What you've built}]
One more word for telling, explaining and describing at B1.
\end{grammarlens}

\subsection*{Your turn}
\begin{itemize}
\raggedright
  \item \emph{Say it:} \emph{torrēns, torrentis}
  \item \emph{Say it:} \emph{torrēns, torrentis}, once more
  \item \emph{Say it:} a sentence of your own with \emph{torrēns, torrentis}
  \item \emph{Recall:} say the Latin for to slide down, then the Latin for a river, then say \emph{torrēns, torrentis} again
\end{itemize}

\begin{tcolorbox}[breakable,colback=teal!4,colframe=teal!35!black,title={Before you move on}]
What does torrēns, torrentis mean? (\textquotedblleft{}A torrent, a rushing stream\textquotedblright{}.) Say it once more.
\end{tcolorbox}
\section[cōnfluēns, cōnfluentis]{cōnfluēns, cōnfluentis — a confluence, the place where two rivers meet}

\label{lesson:LA-C482-confluens}



\begin{quote}
Before the new one: say the Latin for a river, then the Latin for a torrent.
\end{quote}

\subsection*{You'll want to know: cōnfluēns, cōnfluentis}
\textbf{cōnfluēns, cōnfluentis} — \textquotedblleft{}a confluence, the place where two rivers meet\textquotedblright{}.

\begin{quote}
Ad cōnfluentem duōrum flūminum oppidum est.
\end{quote}

\begin{quote}
\emph{At the confluence of two rivers there is a town.}
\end{quote}

Third declension, masculine: \emph{cōnfluēns, cōnfluentis m.}, also plural \emph{cōnfluentēs}. Many towns stand at such a place. English \emph{confluence} comes from the same verb.

\begin{cousinweb}
\emph{cōnfluēns} is the present participle of \emph{cōnfluō} 'to flow together', from \emph{con-} and \emph{fluō} 'to flow' (Lewis \& Short). English \emph{confluence} comes from the same verb.
\end{cousinweb}

\begin{grammarlens}[title={What you've built}]
One more word for telling, explaining and describing at B1.
\end{grammarlens}

\subsection*{Your turn}
\begin{itemize}
\raggedright
  \item \emph{Say it:} \emph{cōnfluēns, cōnfluentis}
  \item \emph{Say it:} \emph{cōnfluēns, cōnfluentis}, once more
  \item \emph{Say it:} a sentence of your own with \emph{cōnfluēns, cōnfluentis}
  \item \emph{Recall:} say the Latin for a river, then the Latin for a torrent, then say \emph{cōnfluēns, cōnfluentis} again
\end{itemize}

\begin{tcolorbox}[breakable,colback=teal!4,colframe=teal!35!black,title={Before you move on}]
What does cōnfluēns, cōnfluentis mean? (\textquotedblleft{}A confluence, the place where two rivers meet\textquotedblright{}.) Say it once more.
\end{tcolorbox}
\section[fluō, fluere]{fluō, fluere — to flow, to run (of water)}

\label{lesson:LA-C482-fluo}



\begin{quote}
Before the new one: say the Latin for a torrent, then the Latin for a confluence.
\end{quote}

\subsection*{You'll want to know: fluō, fluere}
\textbf{fluō, fluere} — \textquotedblleft{}to flow, to run (of water)\textquotedblright{}.

\begin{quote}
Rīvus per prātum lentē fluit.
\end{quote}

\begin{quote}
\emph{A stream flows slowly through the meadow.}
\end{quote}

\emph{fluō, fluere, flūxī, flūxum}. English \emph{fluid}, \emph{fluent} and \emph{flux} come from it.

\begin{cousinweb}
\emph{fluō} is inherited from PIE \emph{b\textsuperscript{h}lew-} 'to swell, to flow' (EtymDB). English \emph{fluid}, \emph{fluent}, \emph{flux} and \emph{influence} come from it (Lewis \& Short).
\end{cousinweb}

\begin{grammarlens}[title={What you've built}]
One more word for telling, explaining and describing at B1.
\end{grammarlens}

\subsection*{Your turn}
\begin{itemize}
\raggedright
  \item \emph{Say it:} \emph{fluō, fluere}
  \item \emph{Say it:} \emph{fluō, fluere}, once more
  \item \emph{Say it:} a sentence of your own with \emph{fluō, fluere}
  \item \emph{Recall:} say the Latin for a torrent, then the Latin for a confluence, then say \emph{fluō, fluere} again
\end{itemize}

\begin{tcolorbox}[breakable,colback=teal!4,colframe=teal!35!black,title={Before you move on}]
What does fluō, fluere mean? (\textquotedblleft{}To flow, to run (of water)\textquotedblright{}.) Say it once more.
\end{tcolorbox}
\section[trānsnatō, trānsnatāre]{trānsnatō, trānsnatāre — to swim across}

\label{lesson:LA-C482-transnato}



\begin{quote}
Before the new one: say the Latin for a confluence, then the Latin for to flow.
\end{quote}

\subsection*{You'll want to know: trānsnatō, trānsnatāre}
\textbf{trānsnatō, trānsnatāre} — \textquotedblleft{}to swim across\textquotedblright{}.

\begin{quote}
Mīlitēs flūmen trānsnatant, quod pōns nūllus est.
\end{quote}

\begin{quote}
\emph{The soldiers swim across the river, because there is no bridge.}
\end{quote}

\emph{trānsnatō, trānsnatāre, trānsnatāvī}, also \emph{trānatō}, with an accusative.

\begin{cousinweb}
\emph{trānsnatō} is built from \emph{trāns} 'across' and \emph{natō} 'to swim' (Lewis \& Short).
\end{cousinweb}

\begin{grammarlens}[title={What you've built}]
One more word for telling, explaining and describing at B1.
\end{grammarlens}

\subsection*{Your turn}
\begin{itemize}
\raggedright
  \item \emph{Say it:} \emph{trānsnatō, trānsnatāre}
  \item \emph{Say it:} \emph{trānsnatō, trānsnatāre}, once more
  \item \emph{Say it:} a sentence of your own with \emph{trānsnatō, trānsnatāre}
  \item \emph{Recall:} say the Latin for a confluence, then the Latin for to flow, then say \emph{trānsnatō, trānsnatāre} again
\end{itemize}

\begin{tcolorbox}[breakable,colback=teal!4,colframe=teal!35!black,title={Before you move on}]
What does trānsnatō, trānsnatāre mean? (\textquotedblleft{}To swim across\textquotedblright{}.) Say it once more.
\end{tcolorbox}
